% TOP_PROBLEM: {"id": 1101211, "title": "Questions in Geometric Group Theory — Q 12.11", "category_id": 17, "category": {"id": 17, "name": "group_theory", "display_name": "Group Theory", "description": "Problems about groups, group actions, representations, and related algebraic structures.", "slug": "group-theory", "order_index": 17, "created_at": "2026-07-31T15:26:25.671Z"}, "status": "partially_solved", "problem_number": "AMR-010-1211", "set_id": 13}
% FIRST_POSTED: 2026-10-02
% ENTRY_KIND: statement_only
% UnsolvedMath ID 1101211; source catalogue code AMR-010-1211
% !TeX program = lualatex
% Definition and sources only; this file is not an LLM solution attempt.
\documentclass[11pt,a4paper]{article}
\usepackage[margin=25mm]{geometry}
\usepackage{fontspec}
\setmainfont{lmroman10-regular.otf}[BoldFont=lmroman10-bold.otf,
 ItalicFont=lmroman10-italic.otf,BoldItalicFont=lmroman10-bolditalic.otf]
\usepackage{amsmath,amssymb,mathtools}
\usepackage{unicode-math}
\setmathfont{latinmodern-math.otf}
\AtBeginDocument{\let\setminus\smallsetminus}
\usepackage{xurl,hyperref}
\hypersetup{colorlinks=true,urlcolor=blue}
\setcounter{secnumdepth}{0}
\setlength{\parindent}{0pt}
\setlength{\parskip}{5pt}
\setlength{\emergencystretch}{3em}
\begin{document}
\section*{Questions in Geometric Group Theory — Q 12.11}\label{1101211}
{\small UnsolvedMath ID 1101211 · AMR-010-1211 · Group Theory · AMR Open Problem Lists}
\subsection{Definitions and mathematical statement}
Let $F_n$, $n\geq2$, be a free group with a fixed free basis. The boundary $\partial F_n$ is the space of infinite reduced words, with two words close when they have a long common initial segment. Equivalently it is the end space of the Cayley tree. Define
\[
 \partial^3 F_n=\{(\xi_1,\xi_2,\xi_3)\in(\partial F_n)^3:
                                      \xi_i\ne\xi_j\ (i\ne j)\}.
\]
Every automorphism induces a boundary homeomorphism and hence acts diagonally on this triple space. Inner automorphisms induce the usual boundary action of $F_n$. Form the quotient
\[
 Y_n=\partial^3F_n/F_n
\]
with its quotient topology. It is a compact totally disconnected space, and the induced action factors through $\operatorname{Out}(F_n)=\operatorname{Aut}(F_n)/\operatorname{Inn}(F_n)$.

Describe the dynamics of this action, particularly the homeomorphism of $Y_n$ induced by a single outer automorphism. Determine its periodic points, orbit closures, recurrent points and invariant closed subsets in terms of algebraic or dynamical information about the automorphism.

The three ends are ordered and pairwise distinct. Passing to the quotient removes only simultaneous translation of the triple; it does not identify arbitrary triples with the same unordered end set. A change of free basis gives an equivalent boundary action, so a proposed dynamical description should not depend on the chosen word coordinates.
\subsection{Short English statement}
Describe how outer automorphisms move triples of free-group ends after simultaneous translation is factored out.
\subsection{Sources}
\begin{itemize}
\item[S1] ULAM.AI and the UnsolvedMath Contributors, supplied problems.json, record 1101211 (AMR-010-1211), AMR Open Problem Lists. Catalogue identity and source transcription; CC BY 4.0.
\url{https://huggingface.co/datasets/ulamai/UnsolvedMath}.
\item[S2] Bestvina - Questions in Geometric Group Theory (pdf) (2004), Question 12.11 (PDF page 21). Original source indexed by AMR-010-1211.
\url{https://www.math.utah.edu/~bestvina/eprints/questions-updated.pdf}.
\end{itemize}
\end{document}
